\scanpage{273}
- \mspage{8} -
\margin{$n^\circ\ 227$}

premier $W$ passant par $P'$, et tel que $z$ ait un zéro sur $W$. Pour le pôle, on considère $z^{-1}$ au lieu de $z$.

Nous allons maintenant étudier le groupe des valeurs associé à un diviseur premier, et montrer qu'il est canoniquement isomorphe aux entiers rationnels, ce qui nous permettra de parler de l'ordre d'un zéro ou d'un pôle d'une fonction.

L'étude étant local, on suppose que $V$ est affine. Considérons d'abord le cas où $V$ est l'espace affine $S^r$. Soit $P=(x_1,\ldots,x_r)$ un point générique de $S^r$. l'anneau $k[x]$ est factoriel, et les diviseurs premiers sont en correspondance biunivoque avec les polynômes irréductibles $p(x)=p$. Chaque fonction $z$ peut s'écrire
\[
z=p^mR(x)
\]
où $R(x)$ est une fonction rationnelle dont le numérateur et le dénominateur sont premiers à $p$. Alors $z$ a un zéro ou un pôle sur $W$ selon que $m>0$ ou $m<0$. On voit que le groupe des valeurs est canoniquement isomorphe aux entiers, et on note
\[
m=\operatorname{ord}_W z.
\]
C'est l'\uline{ordre} de la fonction $z$ sur $W$. Dans le premier cas, $z$ a un zéro d'ordre $m$, et dans le second, un pôle d'ordre $m$.

Passons au cas général, avec un point générique $(x_1,\ldots,x_n)$ de $V$. Soit $(x')$ un point générique d'un diviseur premier $W$ de $V$. On peut supposer $x'_1,\ldots,x'_{r-1}$ algébriquement indépendant. Soit $x_r$ indépendant de $x_1,\ldots,x_{r-1}$. Alors la place qui induit $(x)\longrightarrow(x')$ induit une place sur le sous corps $k(x_1,\ldots,x_r)$ et peut être considérée comme induisant un diviseur premier sur l'espace affine $S^r$. On a vu plus haut que le groupe des valeurs sur ce sous corps est isomorphe aux
\mspage{272}
